\chapter{The Leak}
% OUTLINE Ch6, per Beastmaster's gate re-frame: THE BOUNDARY OF THE
% VARIATIONAL METHOD -- where Ch1-4's clean structures fail (dissipation)
% -- NOT a fifth face, NOT paying the fluid face. The operator's
% viscosity seed carried as labeled intuition wall-to-wall. Laws:
% INTUITION LAW (every intuition labeled; the seed offered AS the
% chapter's motivating intuition, never asserted); CITATION (Rayleigh
% 1870s dissipation function via The Theory of Sound; the NS dissipative
% term named-as-owed; damped-system numerics as bench practice); the
% FENCE first-and-last; drop-to-section fallback was live -- the fence
% holds, the chapter stands. Beastmaster's FINAL arrival-read on the
% fence. Zero NEW device appearances (dose closed at Ch2+Ch4). Gate: Kodo.

The last chapter closed on a noticing: all four faces conserve. The
extremal keeps its stationarity, the invariant its count, the identity
its balance, the bracket its quarry --- nothing leaks. So bring the tour
to an honest bench and let something leak. A pendulum in air: Chapter 1
derives its swing beautifully --- and the swing dies. Each arc is lower
than the last; the energy Chapter 3's summit promised to conserve drains
away into the air's warmth; and after enough of that, the pendulum hangs
at rest with no trace of how it started. The system has \emph{forgotten}.
Where does the variational tour stand before a thing like that? The
honest answer is: awkwardly --- and this chapter is about the
awkwardness, because naming where a method stops is a duty this series
owes its readers. The last volume named where the no-ruler construction
stopped, at the metric; this volume names where the clean variational
structures stop, at the leak.

The awkwardness is precise, not vague, and it has a classical
half-remedy with a name and a date. Friction of the common kind ---
resistance proportional to how fast you push through it --- does not
descend from any potential, so it cannot be folded into the functional
whose variation Chapter 1 taught; the clean Euler--Lagrange house has no
room for it. Rayleigh, in the 1870s, in the work gathered into
\emph{The Theory of Sound}, built the standard annex: a second scalar
--- the dissipation function --- kept alongside the functional, whose
gradient enters the equations of motion as a correction. It works; the
trade uses it daily; and the reader should see it for what it is
architecturally: an \emph{annex}, not a room. The dissipative term rides
beside the variational structure, not inside it. And the deeper faces
lose more than convenience. Noether's bargain --- Chapter 3's summit ---
trades a symmetry of the functional for an exact conservation; the
damped pendulum keeps the symmetry of its laws under time-shift and
conserves nothing, because the bargain's premise --- that the dynamics
descend whole from a functional --- is exactly what the leak breaks. Of
the four faces, only the fourth keeps its full footing: brackets do not
care whether the system remembers, and the numericist integrates damped
systems every working day. The clean faces leak; the squeeze survives;
and that asymmetry is the boundary's shape.

\begin{intuition}
Now the feeling this chapter exists to hand over, at the operator's own
prompting, and the label binds every sentence of it. A system with a
leak \emph{forgets}. Whatever the start --- however violent the launch,
however particular the history --- the transients ring down, the
details drain into the annex, and what remains is the settled thing:
the rest state, the steady hum, the long-run rhythm that no longer
answers questions about the beginning. Now consider what an instrument
sees when it reads such a system. A reading taken after the forgetting
reflects the settled thing only; the itemized bookkeeping the four
clean faces promise --- this much from the launch, this much conserved,
this much exactly cancelled --- is precisely what the leak has
destroyed. The intuition on offer: when a reading resists the clean
accountings --- when it will not itemize, will not conserve, will not
cancel to the coin --- it is \emph{worth asking} whether one is reading
a system mid-forgetting, and the leak is where such readings would
naturally pool. That is a place to look, handed to the reader as a
lens. It is not a finding.
\end{intuition}

Where is the leak itself well understood? The trade's canonical home for
viscosity --- the word for the leak in a flowing medium --- is the
dissipative term of the Navier--Stokes equations, and the reader of this
series knows that address already: it is the fluid face, promised in the
founding text of the series' apparatus, named as owed in every volume
since, and paid in none --- including, emphatically, this one. This
chapter borrows the \emph{word} viscosity for its lens and returns the
word unspent: no fluid is constructed here, no dissipative dynamics
derived, and the fourth station stands exactly as owed as it stood a
chapter ago.

\begin{fence}
The fence, first and last, because this chapter is the volume's nearest
approach to its own failure mode. Everything above the citations in
this chapter is intuition, and labeled; the chapter asserts, as
mathematics, only what Rayleigh's annex and Noether's premise already
contain. The phrase the reader may have heard in this series' wider
conversation --- the viscosity of space-time --- is an \emph{image}, of
the kind the preface licensed: a way of holding the lens, not a claim
about the world, not a property this series has measured, not physics
any volume of it has earned. And the lens itself obeys the lens's law:
that ill-understood readings \emph{might} pool at the leak is a place
to look, offered as such; whether any particular reading is in fact
dissipative must be earned per-reading, off that reading's own
construction, and no reading has been so earned here. The boundary of
the variational method is real, cited, and this chapter's content. The
territory past the boundary is owed --- to the volume that pays the
fluid its debt --- and this book stops at the boundary line, on
purpose, the way its predecessor stopped at a door.
\end{fence}

What remains for this volume is only its accounting: the citation
ledger, chapter by chapter; the intuition ledger the preface promised
--- every picture this book handed over, listed beside the citation it
re-sees; and the honest grades of the faces used. That is the back
matter's business, and the final chapter's.
